\documentclass[12pt,a4paper]{article}
\usepackage[utf8]{inputenc}
\usepackage[spanish]{babel}
\usepackage{geometry}
\geometry{margin=2.5cm}
\usepackage{graphicx}
\usepackage{booktabs}
\usepackage{xcolor}
\usepackage{hyperref}
\usepackage{fancyhdr}
\usepackage{longtable}
\usepackage{array}
\usepackage{colortbl}

\hypersetup{
    colorlinks=true,
    linkcolor=blue,
    urlcolor=blue,
    pdftitle={Informe Benchmark de APIs},
    pdfauthor={Estudiante}
}

\pagestyle{fancy}
\fancyhf{}
\fancyhead[L]{\small Benchmark de APIs — Clean Architecture}
\fancyhead[R]{\small\today}
\fancyfoot[C]{\thepage}
\renewcommand{\headrulewidth}{0.4pt}

\definecolor{headerblue}{RGB}{41,65,114}
\definecolor{lightgray}{RGB}{245,245,245}

\newcolumntype{L}[1]{>{\raggedright\arraybackslash}p{#1}}
\newcolumntype{C}[1]{>{\centering\arraybackslash}p{#1}}
\newcolumntype{R}[1]{>{\raggedleft\arraybackslash}p{#1}}

\begin{document}

% ============================================================
% PORTADA
% ============================================================
\begin{titlepage}
    \centering
    \vspace*{3cm}
    
    {\Huge\bfseries Informe de Resultados}\\[0.5cm]
    {\LARGE Benchmark de APIs con PostgreSQL}\\[1cm]
    {\Large Evaluación Comparativa de 4 Stacks Tecnológicos}\\[2cm]
    
    {\large\textbf{Stacks evaluados:}}\\[0.5cm]
    \begin{tabular}{ll}
        \toprule
        \textbf{Stack} & \textbf{Tecnologías} \\
        \midrule
        Python & Flask + psycopg2 + Gunicorn \\
        Node.js & NestJS + pg (driver nativo) \\
        Java & Spring Boot + JdbcTemplate \\
        Go & net/http + pgx (driver nativo) \\
        \bottomrule
    \end{tabular}
    
    \vfill
    
    {\large Arquitectura: \textbf{Clean Architecture}}\\[0.5cm]
    {\large Herramienta de carga: \textbf{Apache JMeter 5.6.3}}\\[0.5cm]
    {\large Base de datos: \textbf{PostgreSQL 16 (Docker)}}\\[2cm]
    
    {\large \today}
\end{titlepage}

\tableofcontents
\newpage

% ============================================================
% 1. INTRODUCCIÓN Y METODOLOGÍA
% ============================================================
\section{Introducción y Metodología}

El presente informe documenta los resultados de un benchmark comparativo entre cuatro stacks tecnológicos implementando el mismo microservicio: \texttt{POST /api/orders}. Cada stack sigue los principios de \textbf{Clean Architecture}, con separación de capas (Domain, Application, Adapter, Infrastructure) y uso de SQL nativo (sin ORMs).

\subsection{Endpoint evaluado}

\begin{itemize}
    \item \textbf{URL:} \texttt{POST http://localhost:<puerto>/api/orders}
    \item \textbf{Request body:} JSON con \texttt{customerId} y lista de \texttt{items}
    \item \textbf{Response:} JSON con \texttt{orderId}, \texttt{totalAmount}, \texttt{itemsCount}, \texttt{createdAt}
    \item \textbf{Base de datos:} PostgreSQL 16 en contenedor Docker independiente por stack
\end{itemize}

\subsection{Escenarios de carga (JMeter)}

\begin{table}[h]
\centering
\begin{tabular}{lrrr}
\toprule
\textbf{Escenario} & \textbf{Usuarios} & \textbf{Duración (s)} & \textbf{Ramp-up (s)} \\
\midrule
Baja carga   & 10  & 60  & 2  \\
Media carga  & 50  & 120 & 5  \\
Alta carga   & 100 & 180 & 10 \\
Estrés       & 150 & 240 & 15 \\
\bottomrule
\end{tabular}
\caption{Configuración de Thread Groups en JMeter}
\end{table}

\subsection{Métricas registradas}

Por cada escenario se registraron:
\begin{enumerate}
    \item \textbf{Average Response Time} (ms): tiempo promedio de respuesta
    \item \textbf{Percentil 95} (ms): latencia en el percentil 95
    \item \textbf{Throughput} (req/s): solicitudes procesadas por segundo
    \item \textbf{Error rate} (\%): porcentaje de requests fallidos
    \item \textbf{Observaciones}: cuellos de botella, timeouts, comportamiento anómalo
\end{enumerate}

\newpage

% ============================================================
% 2. RESULTADOS INDIVIDUALES
% ============================================================
\section{Resultados Individuales por Stack}

\subsection{Flask (Python)}

\begin{table}[h]
\centering
\small
\begin{tabular}{lrrrr}
\toprule
\textbf{Escenario} & \textbf{Avg (ms)} & \textbf{P95 (ms)} & \textbf{Throughput (req/s)} & \textbf{\% Error} \\
\midrule
Baja carga (10)  & 17,1 & 32 & 563,3 & 0,00 \% \\
Media carga (50) & 86,3 & 146 & 566,3 & 0,00 \% \\
Alta carga (100) & 150,4 & 249 & 646,4 & 0,00 \% \\
Estrés (150)     & 292,9 & 508 & 496,3 & 0,00 \% \\
\bottomrule
\end{tabular}
\caption{Resultados de rendimiento — Flask (Python)}
\end{table}

\textbf{Observaciones:}
\begin{itemize}
    \item El servidor Gunicorn con 4 workers manejó bien las cargas baja y media.
    \item Bajo alta carga (100 usuarios), la latencia promedio aumentó a 150 ms.
    \item En el escenario de estrés (150 usuarios), el throughput disminuyó de 646 a 496 req/s, indicando saturación del pool de workers.
    \item No se registraron errores HTTP ni timeouts en ningún escenario.
    \item El GIL de Python limita el paralelismo real a nivel de proceso; el uso de múltiples workers mitiga pero no elimina el cuello de botella.
\end{itemize}

\subsection{Go (stdlib + pgx)}

\begin{table}[h]
\centering
\small
\begin{tabular}{lrrrr}
\toprule
\textbf{Escenario} & \textbf{Avg (ms)} & \textbf{P95 (ms)} & \textbf{Throughput (req/s)} & \textbf{\% Error} \\
\midrule
Baja carga (10)  & 4,5 & 11 & 2111,4 & 0,00 \% \\
Media carga (50) & 21,3 & 35 & 2274,8 & 0,00 \% \\
Alta carga (100) & 40,6 & 78 & 2384,7 & 0,00 \% \\
Estrés (150)     & 39,7 & 52 & 3624,6 & 0,00 \% \\
\bottomrule
\end{tabular}
\caption{Resultados de rendimiento — Go (stdlib)}
\end{table}

\textbf{Observaciones:}
\begin{itemize}
    \item Go demostró el mejor rendimiento absoluto en todos los escenarios.
    \item La latencia promedio se mantuvo por debajo de 42 ms incluso bajo 150 usuarios concurrentes.
    \item El throughput escaló de 2111 req/s (baja) a 3624 req/s (estrés), mostrando una curva de escalabilidad casi lineal.
    \item No se registraron errores ni fallas de conexión a la base de datos.
    \item El modelo de goroutines + pgxpool permite un manejo eficiente de concurrencia sin el overhead de threads del sistema operativo.
\end{itemize}

\subsection{NestJS (Node.js)}

\textit{Nota: Los benchmarks para este stack no fueron ejecutados en esta sesión. Las tablas deben completarse ejecutando:}

\begin{verbatim}
cd nodejs-nestjs && docker compose up -d --build
./scripts/run-load-test.sh nodejs 5002 low
./scripts/run-load-test.sh nodejs 5002 medium
./scripts/run-load-test.sh nodejs 5002 high
./scripts/run-load-test.sh nodejs 5002 stress
\end{verbatim}

\subsection{Spring Boot (Java)}

\textit{Nota: Los benchmarks para este stack no fueron ejecutados en esta sesión. Las tablas deben completarse ejecutando:}

\begin{verbatim}
cd java-springboot && docker compose up -d --build
./scripts/run-load-test.sh java 5003 low
./scripts/run-load-test.sh java 5003 medium
./scripts/run-load-test.sh java 5003 high
./scripts/run-load-test.sh java 5003 stress
\end{verbatim}

\newpage

% ============================================================
% 3. TABLA COMPARATIVA GRUPAL
% ============================================================
\section{Tabla Comparativa — Entrega Grupal}

\subsection{Comparación de rendimiento por escenario}

\subsubsection{Baja carga (10 usuarios)}

\begin{table}[h]
\centering
\small
\begin{tabular}{lrrrr}
\toprule
\textbf{Stack} & \textbf{Avg (ms)} & \textbf{P95 (ms)} & \textbf{Throughput} & \textbf{\% Error} \\
\midrule
Flask (Python)     & 17,1 & 32 & 563,3 & 0,00 \% \\
Go (stdlib)        & 4,5  & 11 & 2111,4 & 0,00 \% \\
NestJS (Node)      & --   & -- & -- & -- \\
Spring Boot (Java) & --   & -- & -- & -- \\
\bottomrule
\end{tabular}
\caption{Baja carga — Comparativa}
\end{table}

\subsubsection{Media carga (50 usuarios)}

\begin{table}[h]
\centering
\small
\begin{tabular}{lrrrr}
\toprule
\textbf{Stack} & \textbf{Avg (ms)} & \textbf{P95 (ms)} & \textbf{Throughput} & \textbf{\% Error} \\
\midrule
Flask (Python)     & 86,3 & 146 & 566,3 & 0,00 \% \\
Go (stdlib)        & 21,3 & 35 & 2274,8 & 0,00 \% \\
NestJS (Node)      & --   & -- & -- & -- \\
Spring Boot (Java) & --   & -- & -- & -- \\
\bottomrule
\end{tabular}
\caption{Media carga — Comparativa}
\end{table}

\subsubsection{Alta carga (100 usuarios)}

\begin{table}[h]
\centering
\small
\begin{tabular}{lrrrr}
\toprule
\textbf{Stack} & \textbf{Avg (ms)} & \textbf{P95 (ms)} & \textbf{Throughput} & \textbf{\% Error} \\
\midrule
Flask (Python)     & 150,4 & 249 & 646,4 & 0,00 \% \\
Go (stdlib)        & 40,6  & 78  & 2384,7 & 0,00 \% \\
NestJS (Node)      & --    & --  & -- & -- \\
Spring Boot (Java) & --    & --  & -- & -- \\
\bottomrule
\end{tabular}
\caption{Alta carga — Comparativa}
\end{table}

\subsubsection{Estrés (150 usuarios)}

\begin{table}[h]
\centering
\small
\begin{tabular}{lrrrr}
\toprule
\textbf{Stack} & \textbf{Avg (ms)} & \textbf{P95 (ms)} & \textbf{Throughput} & \textbf{\% Error} \\
\midrule
Flask (Python)     & 292,9 & 508 & 496,3 & 0,00 \% \\
Go (stdlib)        & 39,7  & 52  & 3624,6 & 0,00 \% \\
NestJS (Node)      & --    & --  & -- & -- \\
Spring Boot (Java) & --    & --  & -- & -- \\
\bottomrule
\end{tabular}
\caption{Estrés — Comparativa}
\end{table}

\newpage

% ============================================================
% 4. CONSUMO DEL SISTEMA
% ============================================================
\section{Consumo del Sistema (CPU / RAM)}

\textit{Nota: Esta sección debe completarse con datos obtenidos mediante \texttt{docker stats} y \texttt{htop} durante la ejecución de cada prueba.}

\begin{table}[h]
\centering
\small
\begin{tabular}{llrrrrL{4cm}}
\toprule
\textbf{Stack} & \textbf{Escenario} & \textbf{CPU API} & \textbf{RAM API} & \textbf{CPU DB} & \textbf{RAM DB} & \textbf{Observaciones} \\
\midrule
Python & Baja (10)   & -- & -- & -- & -- & (completar) \\
Python & Media (50)  & -- & -- & -- & -- & (completar) \\
Python & Alta (100)  & -- & -- & -- & -- & (completar) \\
Python & Estrés (150)& -- & -- & -- & -- & (completar) \\
\midrule
Go & Baja (10)   & -- & -- & -- & -- & (completar) \\
Go & Media (50)  & -- & -- & -- & -- & (completar) \\
Go & Alta (100)  & -- & -- & -- & -- & (completar) \\
Go & Estrés (150)& -- & -- & -- & -- & (completar) \\
\bottomrule
\end{tabular}
\caption{Consumo de recursos por stack y escenario}
\end{table}

\newpage

% ============================================================
% 5. EVIDENCIAS
% ============================================================
\section{Evidencias}

\textit{Esta sección debe incluir las siguientes capturas de pantalla por cada stack y escenario:}

\begin{enumerate}
    \item \textbf{Summary Report} de JMeter
    \item \textbf{Aggregate Report} de JMeter
    \item \textbf{Pantallazo de JMeter} durante la ejecución de la carga
    \item \textbf{docker ps} mostrando los contenedores activos
    \item \textbf{htop} mostrando el consumo de CPU/RAM del host
    \item \textbf{docker stats} mostrando métricas de los contenedores
\end{enumerate}

\textit{Las capturas deben guardarse en \texttt{docs/evidencias/<stack>/<escenario>/}.}

\newpage

% ============================================================
% 6. CONCLUSIONES
% ============================================================
\section{Conclusiones}

\subsection{Rendimiento}

\begin{enumerate}
    \item \textbf{¿Cuál stack tuvo mejor latencia promedio?}
    \begin{itemize}
        \item Go fue el claro ganador con latencias promedio de 4,5 ms (baja carga) a 40,6 ms (alta carga).
        \item Python (Flask) presentó latencias significativamente mayores, alcanzando 292,9 ms bajo estrés.
    \end{itemize}
    
    \item \textbf{¿Cuál mantuvo el mejor throughput en cargas altas?}
    \begin{itemize}
        \item Go escaló su throughput de 2111 req/s a 3624 req/s bajo estrés.
        \item Python alcanzó un pico de 646 req/s en alta carga pero decayó a 496 req/s bajo estrés, mostrando saturación.
    \end{itemize}
    
    \item \textbf{¿Qué lenguaje fue más estable?}
    \begin{itemize}
        \item Ambos stacks (Python y Go) mantuvieron 0\% de errores en todos los escenarios.
        \item Go mostró mayor estabilidad en latencia, sin degradación abrupta.
    \end{itemize}
\end{enumerate}

\subsection{Eficiencia}

\begin{enumerate}
    \setcounter{enumi}{3}
    \item \textbf{¿Algún stack mostró errores bajo alta carga?}
    \begin{itemize}
        \item No se registraron errores en Python ni Go.
        \item Los assertions de JMeter (status 200 + body con orderId) validaron la integridad de todas las respuestas.
    \end{itemize}
    
    \item \textbf{¿Qué puede explicar las diferencias de rendimiento?}
    \begin{itemize}
        \item \textbf{Go:} modelo de goroutines ligero + pgxpool eficiente permite alta concurrencia sin overhead de threads del SO.
        \item \textbf{Python:} el GIL (Global Interpreter Lock) limita el paralelismo real dentro de un proceso. Aunque Gunicorn usa 4 workers, el overhead de IPC y el modelo WSGI sincrónico generan latencias mayores bajo carga concurrente intensa.
        \item \textbf{Pool de conexiones:} ambos usan pool de 20 conexiones, pero Go reutiliza conexiones de forma más eficiente gracias a su runtime.
    \end{itemize}
\end{enumerate}

\subsection{Recomendación para producción}

\begin{enumerate}
    \setcounter{enumi}{5}
    \item \textbf{¿Cuál stack recomendarían para una empresa que requiere alto desempeño?}
    \begin{itemize}
        \item \textbf{Go} es la recomendación para alto throughput y baja latencia. Su modelo de concurrencia y footprint de memoria reducido lo hacen ideal para microservicios de alta carga.
    \end{itemize}
    
    \item \textbf{¿Cuál stack es más fácil de implementar pero sacrifica rendimiento?}
    \begin{itemize}
        \item \textbf{Python (Flask)} ofrece la curva de aprendizaje más suave y mayor cantidad de recursos educativos, pero su rendimiento es significativamente inferior bajo carga concurrente.
    \end{itemize}
    
    \item \textbf{¿Cuál stack ofrece un mejor balance entre sencillez y velocidad?}
    \begin{itemize}
        \item \textbf{Go} ofrece un excelente balance: sintaxis moderna, compilación rápida, y rendimiento cercano a C/C++.
        \item \textbf{Node.js (NestJS)} podría ser una alternativa intermedia (pendiente de benchmark), ya que su modelo async/IO no bloqueante permite buen rendimiento con curva de aprendizaje moderada.
    \end{itemize}
\end{enumerate}

\newpage

% ============================================================
% 7. ANEXOS
% ============================================================
\section{Anexos}

\subsection{Comandos utilizados}

\begin{verbatim}
# Levantar un stack
cd python-flask && docker compose up -d --build

# Ejecutar benchmark
./scripts/run-load-test.sh python 5001 low

# Monitorear recursos
./scripts/monitor.sh python low

# Extraer métricas automáticamente
./scripts/extract-metrics.sh

# Generar reporte HTML de JMeter
# (automático al ejecutar run-load-test.sh)
\end{verbatim}

\subsection{Estructura del repositorio}

\begin{verbatim}
api-benchmarking-jmeter/
  docs/
    INFORME-AUTO.md          <- Metricas extraidas de CSV
    evidencias/              <- Capturas de pantalla
    INFORME.pdf              <- Este documento
  results/                     <- CSV y reportes HTML de JMeter
  scripts/
    run-load-test.sh         <- Wrapper JMeter
    monitor.sh               <- Docker stats
    extract-metrics.sh       <- Extrae metricas a Markdown
    run-full-benchmark.sh    <- Orquestador completo
  shared/jmeter/               <- Test Plans (.jmx)
  python-flask/
  nodejs-nestjs/
  java-springboot/
  go-stdlib/
\end{verbatim}

\vfill
\begin{center}
\rule{0.5\textwidth}{0.4pt}\\[0.5cm]
\textit{Documento generado automáticamente a partir de los resultados de JMeter.}
\end{center}

\end{document}
